%% Question bank chapter file - generated by bank_tools.py build (do not edit while the bank is still being built:
%% it is regenerated from the per-test data in Question Bank\data). Every question is one bankq block:
%%   \begin{bankq}[id=..., source={...}, paper=.., qno=.., type=..., subtopic={...}, answer={...}, answer_src=key|solved,
%%                 working={...}, difficulty=1-10, diff_parts={K S I C F}, diff_note={...},
%%                 group=<passage id or empty>, also_in={duplicate copies in other tests}]
%% difficulty = min(10, K+S+I+C+F-1): K concepts 1-3, S steps 1-4, I insight 0-2, C computation 0-2, F format 0-1
%%   ...question LaTeX, options, figures (\bfig / \bfigi -> img/)...
%%   \end{bankq}
%% Passages shared by several questions are bankpassage blocks; their questions carry group=<passage id>.
\documentclass[10pt]{article}
\usepackage[a4paper,margin=1.8cm]{geometry}
\usepackage[T1]{fontenc}
\usepackage{amsmath,amssymb,mathtools}
\usepackage{graphicx}
\usepackage[export]{adjustbox}
\usepackage{xcolor,needspace,keyval}
\usepackage{enumitem}
\usepackage{booktabs,array,tabularx}
\usepackage[hidelinks]{hyperref}
\setlength{\parindent}{0pt}
\setlength{\parskip}{3pt}
\makeatletter
\define@key{bq}{id}{\def\bq@id{#1}}\define@key{bq}{source}{\def\bq@src{#1}}\define@key{bq}{paper}{\def\bq@paper{#1}}
\define@key{bq}{qno}{\def\bq@qno{#1}}\define@key{bq}{type}{\def\bq@type{#1}}\define@key{bq}{subtopic}{\def\bq@sub{#1}}
\define@key{bq}{answer}{\def\bq@ans{#1}}\define@key{bq}{answer_src}{\def\bq@asrc{#1}}\define@key{bq}{working}{\def\bq@work{#1}}
\define@key{bq}{group}{\def\bq@group{#1}}\define@key{bq}{also_in}{\def\bq@also{#1}}
\define@key{bq}{difficulty}{\def\bq@diff{#1}}\define@key{bq}{diff_parts}{\def\bq@dparts{#1}}\define@key{bq}{diff_note}{\def\bq@dnote{#1}}
\newenvironment{bankq}[1][]{\def\bq@work{}\def\bq@also{}\def\bq@group{}\def\bq@diff{?}\def\bq@dparts{}\def\bq@dnote{}\setkeys{bq}{#1}%
  \par\needspace{4\baselineskip}\medskip\noindent{\small\textbf{\bq@id}\hfill\color{gray}[\bq@src, Paper \bq@paper, Q\bq@qno]
  \bq@type\ $\cdot$ \bq@sub\ $\cdot$ \textbf{D\bq@diff}}\par\nopagebreak\relax}
  {\par{\small\color{teal}Answer: \textbf{\bq@ans}\ifx\bq@work\@empty\else\ (solved: \bq@work)\fi
  \ifx\bq@also\@empty\else\par\color{gray}Also in: \bq@also\fi}\par}
\define@key{bp}{id}{\def\bp@id{#1}}\define@key{bp}{source}{\def\bp@src{#1}}\define@key{bp}{questions}{\def\bp@qs{#1}}
\newenvironment{bankpassage}[1][]{\setkeys{bp}{#1}\par\needspace{4\baselineskip}\bigskip\noindent
  {\small\textbf{Passage \bp@id}\hfill\color{gray}[\bp@src; questions \bp@qs]}\par\nopagebreak\relax\itshape}{\par}
\makeatother
\newcommand{\bfig}[2]{\par\smallskip{\centering\includegraphics[width=#1pt,max width=0.62\linewidth,max height=0.32\textheight]{img/#2}\par}\smallskip}
\newcommand{\bfigi}[2]{\includegraphics[width=#1pt,max width=\linewidth,max height=0.32\textheight,valign=c]{img/#2}}
\newcommand{\bmissing}[1]{\textit{[figure missing in the source conversion: \detokenize{#1}]}}
\newenvironment{qoptions}{\begin{itemize}[leftmargin=1.8em,itemsep=1pt,topsep=2pt,label={}]}{\end{itemize}}
\newcommand{\opt}[2]{\item[(#1)] #2}
\providecommand{\fwstart}{}\providecommand{\fwend}{}\providecommand{\qfullsplit}{}
\newcolumntype{L}[1]{>{\raggedright\arraybackslash}p{#1}}
\newcolumntype{C}[1]{>{\centering\arraybackslash}p{#1}}
\newcolumntype{R}[1]{>{\raggedleft\arraybackslash}p{#1}}
\providecommand{\sectionhdr}[1]{}\providecommand{\subsectionhdr}[1]{}


\begin{document}
{\Large\bfseries Alkyl Halides --- Question Bank}\par\smallskip
{\small Chemistry chapter C18. Questions from all converted test papers, grouped by source series. Answers: printed key where it matches the paper, otherwise solved (answer\_src=solved, with working). Exact repeats across papers appear once, with also\_in listing the other copies.}\par\medskip\hrule

\section*{TOAS papers}
\begin{bankq}[id=TOAS01-P2-CHEM-S1-Q18, source={TOAS - Test 01 (30 Jan 2022)}, paper=2, qno=18, type=single, subtopic={SN2}, answer={B}, answer_src=key, difficulty=5, diff_parts={K2 S3 I1 C0 F0}, diff_note={Nucleophilicity in protic solvent: charge, size and solvation}]
The order of increasing nucleophilicity of the anions (in polar protic solvent)
\bfig{168.2}{TOAS01_p2_chem_s1_q18_1.png}
\begin{qoptions}
\opt{A}{V $<$ I $<$ III $<$ II $<$ IV}
\opt{B}{III $<$ I $<$ IV $<$ II $<$ V}
\opt{C}{III $<$ II $<$ I $<$ IV $<$ V}
\opt{D}{II $<$ III $<$ I $<$ V $<$ IV}
\end{qoptions}
\end{bankq}
\begin{bankq}[id=TOAS01-P2-CHEM-S1-Q19, source={TOAS - Test 01 (30 Jan 2022)}, paper=2, qno=19, type=single, subtopic={SN2}, answer={D}, answer_src=key, difficulty=3, diff_parts={K1 S3 I0 C0 F0}, diff_note={Leaving group ability order for four sets}]
Which of the following order of leaving group is incorrect ?
\begin{qoptions}
\opt{A}{$\mathrm{F^-<Cl^-<Br^-<I^-}$}
\opt{B}{\bfigi{137.3}{TOAS01_p2_chem_s1_q19_B.png}}
\opt{C}{$\mathrm{CH_3-OH>CH_3-O^-}$}
\opt{D}{$\mathrm{Cl^-<H^-}$}
\end{qoptions}
\end{bankq}
\begin{bankq}[id=TOAS02-P1-CHEM-S1-Q04, source={TOAS - Test 02 (13 Feb 2022)}, paper=1, qno=4, type=single, subtopic={SN2}, answer={B}, answer_src=key, difficulty=2, diff_parts={K1 S2 I0 C0 F0}, diff_note={Four statements on S_N1/S_N2 and nucleophilicity}]
Select correct statement(s) among the following
\begin{qoptions}
\opt{A}{Polar protic solvent favours $\mathrm{S_N2}$ mechanism more than polar aprotic solvent generally}
\opt{B}{Nucleophilic order of halides in polar protic solvent is\\ $\mathrm{I^{\ominus}>Br^{\ominus}>Cl^{\ominus}>F^{\ominus}}$}
\opt{C}{Always perfect racemisation take place in $\mathrm{S_N1}$ mechanism}
\opt{D}{Inversion of configuration in $\mathrm{S_N2}$ mechanism always produces R to S configuration or S to R configuration}
\end{qoptions}
\end{bankq}
\begin{bankq}[id=TOAS04-P1-CHEM-S1-Q03, source={TOAS - Test 04 (15 Mar 2022)}, paper=1, qno=3, type=single, subtopic={SN2}, answer={C}, answer_src=key, difficulty=3, diff_parts={K1 S2 I1 C0 F0}, diff_note={SN2 reactivity: steric and electronic comparison of four halides}]
Which of the following is most reactive towards $\mathrm{S_N2}$ mechanism
\begin{qoptions}
\opt{A}{$\mathrm{CH_3{-}Cl}$}
\opt{B}{$\mathrm{CH_2{=}CH{-}CH_2{-}Cl}$}
\opt{C}{\bfigi{57.5}{TOAS04_p1_chem_s1_q03_C.png}}
\opt{D}{\bfigi{51.7}{TOAS04_p1_chem_s1_q03_D.png}}
\end{qoptions}
\end{bankq}
\begin{bankpassage}[id=TOAS04-P2-CHEM-S2-G9-10, source={TOAS - Test 04 (15 Mar 2022)}, questions={9-10}]
Alkyl halides can undergo $\mathrm{S_N^{\,1}}$, $\mathrm{S_N^{\,2}}$, $\mathrm{E_1}$, $\mathrm{E_2}$, $\mathrm{E_{1cb}}$ depending on various conditions.
\bfig{128.3}{TOAS04_p2_chem_s2_para9_10_1.png}
\end{bankpassage}
\begin{bankq}[id=TOAS04-P2-CHEM-S2-Q09, source={TOAS - Test 04 (15 Mar 2022)}, paper=2, qno=9, type=single, subtopic={E2}, answer={C}, answer_src=key, difficulty=5, diff_parts={K2 S3 I1 C0 F0}, diff_note={Substrate and base/solvent decide substitution or elimination for product Y}, group=TOAS04-P2-CHEM-S2-G9-10]
Identify `Y' :
\begin{qoptions}
\opt{A}{\bfigi{86.2}{TOAS04_p2_chem_s2_q09_A.png}}
\opt{B}{\bfigi{69.4}{TOAS04_p2_chem_s2_q09_B.png}}
\opt{C}{\bfigi{66.8}{TOAS04_p2_chem_s2_q09_C.png}}
\opt{D}{Both (B) \& (C)}
\end{qoptions}
\end{bankq}
\begin{bankq}[id=TOAS04-P2-CHEM-S2-Q10, source={TOAS - Test 04 (15 Mar 2022)}, paper=2, qno=10, type=single, subtopic={E2}, answer={D}, answer_src=key, difficulty=5, diff_parts={K2 S3 I1 C0 F0}, diff_note={Base strength and leaving group decide E2 and the orientation of the alkene}, group=TOAS04-P2-CHEM-S2-G9-10]
\bfig{206.5}{TOAS04_p2_chem_s2_q09_4.png}
Select the correct statement :
\begin{qoptions}
\opt{A}{Formation of U follows $\mathrm{E^2}$ mechanism to give saytzeff product}
\opt{B}{Formation of U follows $\mathrm{E_{1cb}}$ mechanism to give saytzeff product}
\opt{C}{Formation of U follows $\mathrm{E^1}$ mechanism to give Hoffmann product}
\opt{D}{Formation of U follows $\mathrm{E^2}$ mechanism to give Hoffmann product}
\end{qoptions}
\end{bankq}
\begin{bankq}[id=TOAS05-P1-CHEM-S1-Q02, source={TOAS - Test 05 (03 Apr 2022)}, paper=1, qno=2, type=multi, subtopic={SN1}, answer={A,B,D}, answer_src=key, difficulty=7, diff_parts={K2 S4 I1 C0 F1}, diff_note={SN1, oxymercuration, SN2 and acid elimination: which give a pair of diastereomers}]
Which of the following reaction result the formation of a pair of diastereomers?
\begin{qoptions}
\opt{A}{\bfigi{90.6}{TOAS05_p1_chem_s1_q02_A.png}}
\opt{B}{\bfigi{78.5}{TOAS05_p1_chem_s1_q02_B.png}}
\opt{C}{\bfigi{90.3}{TOAS05_p1_chem_s1_q02_C.png}}
\opt{D}{\bfigi{29.7}{TOAS05_p1_chem_s1_q02_D.png}}
\end{qoptions}
\end{bankq}
\begin{bankq}[id=TOAS05-P1-CHEM-S3-Q02, source={TOAS - Test 05 (03 Apr 2022)}, paper=1, qno=2, type=numerical, subtopic={SN2}, answer={5.00}, answer_src=key, difficulty=6, diff_parts={K1 S4 I1 C0 F1}, diff_note={SN2 reactivity compared with isopropyl chloride for nine halides}]
Number of compounds which are more reactive than $\mathrm{CH_3{-}CH(CH_3){-}Cl}$ by $\mathrm{S_N^{\,2}}$ mechanism?
\bfig{178.8}{TOAS05_p1_chem_s3_q02_2.png}
\end{bankq}
\begin{bankq}[id=TOAS05-P1-CHEM-S3-Q03, source={TOAS - Test 05 (03 Apr 2022)}, paper=1, qno=3, type=numerical, subtopic={E2}, answer={3.00}, answer_src=key, difficulty=6, diff_parts={K1 S4 I1 C0 F1}, diff_note={Which substrates have a beta hydrogen and a good leaving group for E2}]
How many of the following substrate can give E2 elimination reaction on treatment with base / $\Delta$
\bfig{243.5}{TOAS05_p1_chem_s3_q03_1.png}
\end{bankq}
\begin{bankpassage}[id=TOAS07-P2-CHEM-S2-G7-8, source={TOAS - Test 07 (01 May 2022)}, questions={7-8}]
In an E2 reaction, following one step mechanism is involved.
\bfig{292.9}{TOAS07_p2_chem_s2_para7_8_1.png}

As shown in the above mechanism, a $\beta$-proton anti to leaving group is abstracted by the alkoxide base. If a less bulky base is used, $\beta$-proton is eliminated giving most substituted, most stable product as the major product. On the otherhand if a bulky base is used, $\beta$-proton is abstracted from $\beta$-carbon where there is least steric hindrance and this gives least substituted, although least stable, alkene as the major product.
\end{bankpassage}
\begin{bankq}[id=TOAS07-P2-CHEM-S2-Q07, source={TOAS - Test 07 (01 May 2022)}, paper=2, qno=7, type=single, subtopic={E2}, answer={C}, answer_src=key, difficulty=5, diff_parts={K2 S3 I1 C0 F0}, diff_note={Bulky vs small base fixes Hofmann or Zaitsev elimination product}, group=TOAS07-P2-CHEM-S2-G7-8]
What is the major product in the following elimination reaction ?
\bfig{113.4}{TOAS07_p2_chem_s2_q07_1.png}
\begin{qoptions}
\opt{A}{\bfigi{48.5}{TOAS07_p2_chem_s2_q07_A.png}}
\opt{B}{\bfigi{50.6}{TOAS07_p2_chem_s2_q07_B.png}}
\opt{C}{\bfigi{48.7}{TOAS07_p2_chem_s2_q07_C.png}}
\opt{D}{\bfigi{51.9}{TOAS07_p2_chem_s2_q07_D.png}}
\end{qoptions}
\end{bankq}
\begin{bankq}[id=TOAS07-P2-CHEM-S2-Q08, source={TOAS - Test 07 (01 May 2022)}, paper=2, qno=8, type=single, subtopic={E2}, answer={D}, answer_src=key, difficulty=4, diff_parts={K1 S3 I1 C0 F0}, diff_note={Count the possible E2 alkenes with anti beta hydrogens}, group=TOAS07-P2-CHEM-S2-G7-8]
In the following reaction, how many elimination products would be formed in principle by E2 mechanism ?
\bfig{143.4}{TOAS07_p2_chem_s2_q08_1.png}
\begin{qoptions}
\opt{A}{2}
\opt{B}{3}
\opt{C}{4}
\opt{D}{5}
\end{qoptions}
\end{bankq}
\begin{bankpassage}[id=TOAS08-P1-CHEM-S2-G9-10, source={TOAS - Test 08 (15 May 2022)}, questions={9-10}]
Alkyl halides can undergo $\mathrm{S_N1}$, $\mathrm{S_N2}$, $\mathrm{E_1}$, $\mathrm{E_2}$, $\mathrm{E_{1cb}}$ depending on various conditions.
\bfig{159.0}{TOAS08_p1_chem_s2_para9_10_1.png}
\end{bankpassage}
\begin{bankq}[id=TOAS08-P1-CHEM-S2-Q09, source={TOAS - Test 08 (15 May 2022)}, paper=1, qno=9, type=single, subtopic={E2}, answer={D}, answer_src=key, difficulty=5, diff_parts={K2 S3 I1 C0 F0}, diff_note={Elimination mechanism and product type for U}, group=TOAS08-P1-CHEM-S2-G9-10]
Select the correct statement :
\begin{qoptions}
\opt{A}{Formation of $U$ follows $\mathrm{E^2}$ mechanism to give saytzeff product}
\opt{B}{Formation of $U$ follows $\mathrm{E_{1cb}}$ mechanism to give saytzeff product}
\opt{C}{Formation of $U$ follows $\mathrm{E^1}$ mechanism to give Hoffmann product}
\opt{D}{Formation of $U$ follows $\mathrm{E^2}$ mechanism to give Hoffmann product}
\end{qoptions}
\end{bankq}
\begin{bankq}[id=TOAS08-P1-CHEM-S2-Q10, source={TOAS - Test 08 (15 May 2022)}, paper=1, qno=10, type=single, subtopic={E2}, answer={C}, answer_src=key, difficulty=5, diff_parts={K2 S3 I1 C0 F0}, diff_note={Product Y of the elimination/addition sequence}, group=TOAS08-P1-CHEM-S2-G9-10]
Identify `$Y$' :
\begin{qoptions}
\opt{A}{\bfigi{76.4}{TOAS08_p1_chem_s2_q10_A.png}}
\opt{B}{\bfigi{61.8}{TOAS08_p1_chem_s2_q10_B.png}}
\opt{C}{\bfigi{59.3}{TOAS08_p1_chem_s2_q10_C.png}}
\opt{D}{Both (B) \& (C)}
\end{qoptions}
\end{bankq}
\begin{bankq}[id=TOAS08-P2-CHEM-S1-Q05, source={TOAS - Test 08 (15 May 2022)}, paper=2, qno=5, type=multi, subtopic={SN2}, answer={A,B,D}, answer_src=key, difficulty=6, diff_parts={K1 S4 I1 C0 F1}, diff_note={Which reactions are bimolecular (rate depends on nucleophile)}]
Which of the following reaction will go faster if the concentration of the nucleophile / base is increased?
\begin{qoptions}
\opt{A}{\bfigi{113.1}{TOAS08_p2_chem_s1_q05_A.png}}
\opt{B}{\bfigi{134.5}{TOAS08_p2_chem_s1_q05_B.png}}
\opt{C}{$\mathrm{(CH_3)_3CCl}\ \xrightarrow[\mathrm{CH_3COOH}]{\mathrm{H_2O}}$}
\opt{D}{\bfigi{111.0}{TOAS08_p2_chem_s1_q05_D.png}}
\end{qoptions}
\end{bankq}
\end{document}
